%% LyX 2.4.4 created this file.  For more info, see https://www.lyx.org/.
%% Do not edit unless you really know what you are doing.
\documentclass[english]{article}
\usepackage[T1]{fontenc}
\usepackage[latin9]{inputenc}
\usepackage{units}
\usepackage{enumitem}
\usepackage{amsmath}
\usepackage{amsthm}
\usepackage{amssymb}
\usepackage{geometry}
\geometry{verbose,tmargin=1in,bmargin=1in,lmargin=1in,rmargin=1in}

\makeatletter
%%%%%%%%%%%%%%%%%%%%%%%%%%%%%% Textclass specific LaTeX commands.
\numberwithin{equation}{section}
\numberwithin{figure}{section}
\newlength{\lyxlabelwidth}      % auxiliary length 

%%%%%%%%%%%%%%%%%%%%%%%%%%%%%% User specified LaTeX commands.
\usepackage{fancyhdr}
\usepackage{lastpage}
\pagestyle{fancy}
\usepackage{enumitem}
\usepackage{ifthen}
\everymath=\expandafter{\the\everymath\displaystyle}
%\DeclareMathSizes{10}{10}{10}{10}
\usepackage{multicol}

\usepackage{tikz}
\usepackage{tikz-3dplot}
\usetikzlibrary{calc,arrows}

\usetikzlibrary{patterns}
\usetikzlibrary{plotmarks}
\usepackage{pgfplots}
\pgfplotsset{compat=newest}

\usepackage{calc}
\usepackage[nomessages]{fp}

\usepackage{datetime}
% Added by lyx2lyx
\setlength{\parskip}{\smallskipamount}
\setlength{\parindent}{0pt}

\makeatother

\usepackage{babel}
\begin{document}
\renewcommand{\author}{Dr.\,James Ashe}

\newcommand{\classNum}{336}

\newcommand{\classSec}{A}

\newcommand{\examType}{Exam 1}

\renewcommand{\date}{\formatdate{20}{9}{2013}}

%%First page's header%%%%%%%%%%%%%%%%%%%%%%
\fancypagestyle{firststyle} %%header settings for first pagr
{
	\fancyhead[L]{MTH \classNum{}(\classSec{}) \author{}\\
	\textbf{\examType{}} ~\date{}}
	\fancyhead[C]{}
	\fancyhead[R]{\textbf{Solutions}}
	\setlength{\headsep}{14pt}
}
%%%%%%%%%%%%%%%%%%%%%%%%%%%%%%%%%%%%%%%%%%

%%Next page's header%%%%%%%%%%%%%%%%%%%%%%%
\setlist{nolistsep}
\lhead{\classNum{}(\classSec{})} 
\chead{\textbf{Solutions}} 
\cfoot{\thepage{}~of~\pageref{LastPage}} 
\setlength{\headsep}{5pt} 
%%%%%%%%%%%%%%%%%%%%%%%%%%%%%%%%%%%%%%%%%%%

%%Various settings; see the preamble for other settings%%%%%
%\setenumerate[0]{leftmargin=12pt,itemindent=0pt,labelwidth=0pt}
\setlist{topsep=.5pt}
\renewcommand{\labelenumi}{\textbf{\arabic{enumi}.}}
\renewcommand{\labelenumii}{\textbf{\alph{enumii}.}}
\thispagestyle{firststyle}
%%%%%%%%%%%%%%%%%%%%%%%%%%%%%%%%%%%%%%%%%%

\newcommand{\xmax}{0}
\newcommand{\xmin}{-\xmax}
\newcommand{\xstep}{0}
\newcommand{\ymax}{0}
\newcommand{\ymin}{-\ymax}
\newcommand{\ystep}{0} 
\newcommand{\dspace}{\vspace{1cm}}
\newcommand{\xa}{0}
\newcommand{\xb}{0}
\newcommand{\ya}{1}
\newcommand{\yb}{1}

\newcommand{\xspace}{\vfill}
\newcommand{\width}{4cm}
\newcommand{\height}{4cm}
\renewcommand{\labelenumiii}{}Show your work. 
\begin{enumerate}
\item {[}3 pts each{]} Let $\mathbf{x}=(1,2,2)$ and $\mathbf{y}=(3,0,2)$.

\begin{enumerate}
\item Find $\mathbf{x}-\mathbf{y}$.

\begin{enumerate}
\item $=(1,2,2)-(3,0,2)=(1-3,2-0,2-2)=(-2,2,0)$.
\end{enumerate}
\item Find $\mathbf{x\cdot y}$.

\begin{enumerate}
\item $=(1,2,2)\bullet(3,0,2)=(1\cdot3,2\cdot0,2\cdot2)=(3,0,4)$
\end{enumerate}
\item Find $\Vert\mathbf{y}\Vert$.

\begin{enumerate}
\item $=\sqrt{3^{2}+0^{2}+2^{2}}=\sqrt{13}$
\end{enumerate}
\item Find a unit vector in the direction of $y$. 

\begin{enumerate}
\item A unit vector has magnitude $1$. So we need to scale $\mathbf{y}$
by its own magnitude. $\nicefrac{\mathbf{y}}{\Vert\mathbf{y}\Vert}=\nicefrac{(3,0,2)}{\sqrt{13}}=(\nicefrac{3}{\sqrt{13}},0,\nicefrac{2}{\sqrt{13}})$.\xspace
\end{enumerate}
\end{enumerate}
\end{enumerate}
\begin{enumerate}[resume]
\item {[}3 pts each{]} Use the provided graph to sketch the following vectors.\begin{multicols}{3}

\begin{enumerate}
\item $-\frac{1}{2}\mathbf{y}$\\
\renewcommand{\xa}{1}
\renewcommand{\xb}{3}
\renewcommand{\ya}{4}
\renewcommand{\yb}{-1}\renewcommand{\xmax}{5}
\renewcommand{\xmin}{-5}
\renewcommand{\xstep}{0} %must be xmin+n
\renewcommand{\ymax}{5}
\renewcommand{\ymin}{-5}
\renewcommand{\ystep}{0} %must be ymin+n
%%%%%%%
\begin{tikzpicture} 
\begin{axis}[height=\height, 
	width=\width, 
	axis lines=middle,
	%minor x tick num={0},
	%minor y tick num={0},
	%grid=both,
	%xtick={0,50,...,250},
	xtick=\empty,
	%ytick={0,10,20,...,40},
	ytick=\empty,
	%axis line style={-latex}, 
	xmin=\xmin,xmax=\xmax, 
	ymin=\ymin,ymax=\ymax,
	%restrict y to domain=-1:10,
	%no markers, 
	xlabel={$$},ylabel={$$}, 
	%every axis x label/.append style={anchor=north}, 
	%every axis y label/.append style={anchor=east}
	]
	
	\draw[blue,thick,->] (axis cs:0,0) -- (axis cs:\xa,\xb) node [right] {x};
	\draw[blue,thick,->] (axis cs:0,0) -- (axis cs:\ya,\yb) node [below] {y};
	
	\draw[red, very thick,->] (axis cs:0,0) -- (axis cs:-2,.5);
	
	%%\addplot[color=red,solid,mark=*] coordinates {(-3,-4)};
	%%\addplot[domain=\xmin:\xmax,thick,samples=100,draw=red,<->]{-(x-2)^2+5} [yshift=8pt] node[pos=0.7,right] {$g(x)$};
	
\end{axis}
\end{tikzpicture}
\end{enumerate}
\begin{enumerate}[resume]
\item $\mathbf{x}+\mathbf{y}$\\
\renewcommand{\xa}{1}
\renewcommand{\xb}{3}
\renewcommand{\ya}{4}
\renewcommand{\yb}{-1}\renewcommand{\xmax}{5}
\renewcommand{\xmin}{-5}
\renewcommand{\xstep}{0} %must be xmin+n
\renewcommand{\ymax}{5}
\renewcommand{\ymin}{-5}
\renewcommand{\ystep}{0} %must be ymin+n
%%%%%%%
\begin{tikzpicture} 
\begin{axis}[height=\height, 
	width=\width, 
	axis lines=middle,
	%minor x tick num={0},
	%minor y tick num={0},
	%grid=both,
	%xtick={0,50,...,250},
	xtick=\empty,
	%ytick={0,10,20,...,40},
	ytick=\empty,
	%axis line style={-latex}, 
	xmin=\xmin,xmax=\xmax, 
	ymin=\ymin,ymax=\ymax,
	%restrict y to domain=-1:10,
	%no markers, 
	xlabel={$$},ylabel={$$}, 
	%every axis x label/.append style={anchor=north}, 
	%every axis y label/.append style={anchor=east}
	]
	
	\draw[blue,thick,->] (axis cs:0,0) -- (axis cs:\xa,\xb) node [right] {x};
	\draw[blue,thick,->] (axis cs:0,0) -- (axis cs:\ya,\yb) node [below] {y};
	
	\draw[red,very thick,->] (axis cs:0,0) -- (axis cs:5,2);
	\draw[blue,dashed,->] (axis cs:1,3) -- (axis cs:5,2);
	\draw[blue,dashed,->] (axis cs:4,-1) -- (axis cs:5,2);
	
	
	%%\addplot[color=red,solid,mark=*] coordinates {(-3,-4)};
	%%\addplot[domain=\xmin:\xmax,thick,samples=100,draw=red,<->]{-(x-2)^2+5} [yshift=8pt] node[pos=0.7,right] {$g(x)$};
	
\end{axis}
\end{tikzpicture}
\item $\mathbf{x}-\mathbf{y}$\\
\renewcommand{\xa}{1}
\renewcommand{\xb}{3}
\renewcommand{\ya}{4}
\renewcommand{\yb}{-1}\renewcommand{\xmax}{5}
\renewcommand{\xmin}{-5}
\renewcommand{\xstep}{0} %must be xmin+n
\renewcommand{\ymax}{5}
\renewcommand{\ymin}{-5}
\renewcommand{\ystep}{0} %must be ymin+n
%%%%%%%
\begin{tikzpicture} 
\begin{axis}[height=\height, 
	width=\width, 
	axis lines=middle,
	%minor x tick num={0},
	%minor y tick num={0},
	%grid=both,
	%xtick={0,50,...,250},
	xtick=\empty,
	%ytick={0,10,20,...,40},
	ytick=\empty,
	%axis line style={-latex}, 
	xmin=\xmin,xmax=\xmax, 
	ymin=\ymin,ymax=\ymax,
	%restrict y to domain=-1:10,
	%no markers, 
	xlabel={$$},ylabel={$$}, 
	%every axis x label/.append style={anchor=north}, 
	%every axis y label/.append style={anchor=east}
	]
	
	\draw[blue,thick,->] (axis cs:0,0) -- (axis cs:\xa,\xb) node [right] {x};
	\draw[blue,thick,->] (axis cs:0,0) -- (axis cs:\ya,\yb) node [below] {y};
	
	\draw[red,very thick,->] (axis cs:4,-1) -- (axis cs:1,3);
	
	%%\addplot[color=red,solid,mark=*] coordinates {(-3,-4)};
	%%\addplot[domain=\xmin:\xmax,thick,samples=100,draw=red,<->]{-(x-2)^2+5} [yshift=8pt] node[pos=0.7,right] {$g(x)$};
	
\end{axis}
\end{tikzpicture}
\end{enumerate}
\end{multicols}
\item {[}5 pts each{]} Give Cartesian equations of the given hyperplanes.

\begin{enumerate}
\item The hyperplane in $\mathbb{R^{\mbox{2}}}$ passing through $(1,1)$
and orthogonal to the line $\mathbf{x}=(0,2)+t(2,1)$

\begin{enumerate}
\item []It is orthogonal to the given line so $\mathbf{a}=(2,1)$ is a
normal vector. 
\begin{eqnarray*}
\mathbf{a}\cdot\mathbf{x} & = & \mathbf{a}\cdot\mathbf{x}_{0}\\
(2,1)\cdot(x_{1},x_{2}) & = & (2,1)\cdot(1,1)\\
2x_{1}+x_{2} & = & 3
\end{eqnarray*}
\end{enumerate}
\end{enumerate}
\begin{enumerate}[resume]
\item The hyperplane in $\mathbb{R}^{\mbox{3}}$ which contains the point
$(0,1,0)$ and has normal vector $(1,2,1)$.

\begin{enumerate}[resume]
\item []$\mathbf{a}=(1,2,1)$
\begin{eqnarray*}
\mathbf{a}\cdot\mathbf{x} & = & \mathbf{a}\cdot\mathbf{x}_{0}\\
(1,2,1)\cdot(x_{1},x_{2},x_{3}) & = & (1,2,1)\cdot(0,1,0)\\
x_{1}+2x_{2}+x_{3} & = & 2
\end{eqnarray*}
\end{enumerate}
\end{enumerate}
\item Consider the matrix,
\[
A=\begin{bmatrix}\begin{array}{rrrr}
0 & 1 & 0 & 2\\
2 & 0 & 0 & 2\\
0 & 1 & 2 & 1
\end{array}\end{bmatrix}
\]

\begin{enumerate}
\item {[}4 pts{]} Find the reduced echelon form of $A$.

\begin{enumerate}
\item []
\[
A=\begin{bmatrix}\begin{array}{rrrr}
0 & 1 & 0 & 2\\
2 & 0 & 0 & 2\\
0 & 1 & 2 & 1
\end{array}\end{bmatrix}\rightsquigarrow\begin{array}{c}
\nicefrac{1}{2}\mbox{R2}\\
\mbox{R1}\\
\mbox{R3}
\end{array}\begin{bmatrix}\begin{array}{rrrr}
1 & 0 & 0 & 1\\
0 & 1 & 0 & 2\\
0 & 1 & 2 & 1
\end{array}\end{bmatrix}\rightsquigarrow\begin{array}{c}
\mbox{}\\
\\\nicefrac{1}{2}(\mbox{R3-R2})
\end{array}\begin{bmatrix}\begin{array}{rrrr}
1 & 0 & 0 & 1\\
0 & 1 & 0 & 2\\
0 & 0 & 1 & -\nicefrac{1}{2}
\end{array}\end{bmatrix}
\]
\end{enumerate}
\end{enumerate}
\begin{enumerate}[resume]
\item {[}4pts{]} Give the general solution of $A\mathbf{x}=\mathbf{0}$
in \emph{standard }form.

\begin{enumerate}[resume]
\item []$x_{1}+x_{4}=0$, $x_{2}+2x_{4}=0$, and $x_{3}-\nicefrac{1}{2}x_{4}=0$.
So $\mathbf{x}=(-x_{4},-2x_{4},\nicefrac{1}{2}x_{4},x_{4})=x_{4}(1,2-\nicefrac{1}{2},1)$
or 
\[
\mathbf{x}=x_{4}\begin{bmatrix}\begin{array}{r}
-1\\
-2\\
\nicefrac{1}{2}\\
1
\end{array}\end{bmatrix}
\]
\end{enumerate}
\item {[}2 pts{]} Is the reduced echelon form of any matrix unique? i.e.,
if two different people compute the reduced echelon form of the same
matrix, will they always arrive at the exact same result?

\begin{enumerate}[resume]
\item []\emph{YES.}
\end{enumerate}
\end{enumerate}
\item Consider the system of equations 
\[
\begin{array}{rrrr}
x_{1} &  & +x_{3} & =3\\
 & x_{2} & +2x_{3} & =1\\
 & 2x_{2} & +4x_{3} & =2
\end{array}
\]

\begin{enumerate}
\item {[}2 pts{]} Express the system as an augmented matrix.

\begin{enumerate}
\item []
\[
[A\mathbf{x}|\mathbf{b}]=\begin{bmatrix}\begin{array}{rrr|r}
1 & 0 & 1 & 3\\
0 & 1 & 2 & 1\\
0 & 2 & 4 & 2
\end{array}\end{bmatrix}
\]
\end{enumerate}
\item {[}4 pts{]} Find the reduced echelon form of this augmented matrix.
\[
\begin{bmatrix}\begin{array}{rrr|r}
1 & 0 & 1 & 3\\
0 & 1 & 2 & 1\\
0 & 2 & 4 & 2
\end{array}\end{bmatrix}\rightsquigarrow\begin{array}{c}
\\\\\mbox{R3}-2\mbox{R2}
\end{array}\begin{bmatrix}\begin{array}{rrr|r}
1 & 0 & 1 & 3\\
0 & 1 & 2 & 1\\
0 & 0 & 0 & 0
\end{array}\end{bmatrix}
\]
\item {[}4 pts{]} Give the general solution to the system of equations.

\begin{enumerate}
\item []$x_{1}+x_{3}=3$ and $x_{2}+2x_{3}=1$. So then $\mathbf{x}=(3-x_{3},1-2x_{3},x_{3})=(3,1,0)+x_{3}(-1,-2,1)$
or 
\[
\mathbf{x}=\begin{bmatrix}\begin{array}{r}
3\\
1\\
0
\end{array}\end{bmatrix}+x_{3}\begin{bmatrix}\begin{array}{r}
-1\\
-2\\
1
\end{array}\end{bmatrix}
\]
\end{enumerate}
\end{enumerate}
\item {[}3 pts each{]} Consider the matrices 
\[
A=\begin{bmatrix}2 & 0 & -2\\
0 & 0 & 0\\
0 & 1 & 0
\end{bmatrix}\mbox{ and }\mathbf{b}=\begin{bmatrix}\begin{array}{r}
1\\
1\\
1
\end{array}\end{bmatrix}
\]

\begin{enumerate}
\item How many solutions does the system of equations represented by $A\mathbf{x}=\mathbf{b}$
have?

\begin{enumerate}
\item []None. 
\end{enumerate}
\end{enumerate}
\begin{enumerate}[resume]
\item What is the \emph{rank }of $A$?

\begin{enumerate}[resume]
\item []It has one row of zeros. So the $\mbox{rank}(A)=2$ (number of
rows minus 1).
\end{enumerate}
\item How many constraints does the solutions of $A\mathbf{x}=\mathbf{b_{1}}$
have if 
\[
\mathbf{b_{1}}=\begin{bmatrix}\begin{array}{r}
b_{1}\\
b_{2}\\
b_{3}
\end{array}\end{bmatrix}
\]

\begin{enumerate}[resume]
\item []It will have 1 constraint equation since it has 1 row of zeros. 
\end{enumerate}
\end{enumerate}
\item Complete the following

\begin{enumerate}
\item {[}2 pts{]} Suppose that the vectors $\mathbf{v_{1}}$, $\mathbf{v_{2}}$,
and $\mathbf{v_{3}}$ span the hyperplane $\mathcal{H}$, and $\mathbf{w}=2\mathbf{v_{1}}-\mathbf{v_{2}}+3\mathbf{v_{3}}$.
Is $\mathbf{w}$ in $\mathcal{H}$? 

\begin{enumerate}
\item []YES. Since $\mathbf{w}$ is a linear combination of $\mathbf{v_{1}}$,
$\mathbf{v_{2}}$, and $\mathbf{v_{3}}$. 
\end{enumerate}
\end{enumerate}
\begin{enumerate}[resume]
\item {[}2 pts{]} Suppose that the $n\times n$ matrix $A$ is nonsingular.
Does $A\mathbf{x}=\mathbf{b}$ have a solution? If so, is it unique?
(Recall that nonsingular means that the reduced echelon form of the
matrix does not have a row with all zeros.) 

\begin{enumerate}[resume]
\item []YES it has a solution and YES it is unique. If $A$ is nonsingular
then its reduced echelon form will have a diagonal of $1$'s with
zeros elsewhere. So the augmented matrix will look like, 
\[
\begin{bmatrix}\begin{array}{rrr|r}
1 & 0 & 0 & b_{1}\\
\vdots & \ddots & \vdots & \vdots\\
0 & 0 & 1 & b_{n}
\end{array}\end{bmatrix}
\]
So the solution will be $x_{1}=b_{1}$, $x_{2}=b_{2}$, $\dots,x_{n}=b_{n}$.
\end{enumerate}
\item {[}2 pts{]} Let $\mathbf{x}$ and $\mathbf{y}$ be vectors shown in
the figure below. Does $\mathbf{x}=\mathbf{y}$?\\
\renewcommand{\xa}{1}
\renewcommand{\xb}{3}
\renewcommand{\ya}{3}
\renewcommand{\yb}{5}\renewcommand{\xmax}{6}
\renewcommand{\xmin}{-1}
\renewcommand{\xstep}{0} %must be xmin+n
\renewcommand{\ymax}{6}
\renewcommand{\ymin}{-1}
\renewcommand{\ystep}{0} %must be ymin+n
%%%%%%%
\begin{tikzpicture} 
\begin{axis}[height=\height, 
	width=\width, 
	axis lines=middle,
	%minor x tick num={0},
	%minor y tick num={0},
	%grid=both,
	%xtick={0,50,...,250},
	xtick=\empty,
	%ytick={0,10,20,...,40},
	ytick=\empty,
	%axis line style={-latex}, 
	xmin=\xmin,xmax=\xmax, 
	ymin=\ymin,ymax=\ymax,
	%restrict y to domain=-1:10,
	%no markers, 
	xlabel={$$},ylabel={$$}, 
	%every axis x label/.append style={anchor=north}, 
	%every axis y label/.append style={anchor=east}
	]
	
	%\draw[red,thick,<-] (axis cs:0,0) -- (axis cs:\ya,\yb);
	\draw[red,thick,->] (axis cs:0,0) -- (axis cs:\xa,\xb) node [right] {x};
	\draw[blue,thick,->] (axis cs:2,2) -- (axis cs:\ya,\yb) node [right] {y};
	
	%%\addplot[color=red,solid,mark=*] coordinates {(-3,-4)};
	%%\addplot[domain=\xmin:\xmax,thick,samples=100,draw=red,<->]{-(x-2)^2+5} [yshift=8pt] node[pos=0.7,right] {$g(x)$};
	
\end{axis}
\end{tikzpicture}

\begin{enumerate}[resume]
\item []YES.
\end{enumerate}
\item {[}4 pts{]} Let 
\[
A=\begin{bmatrix}\begin{array}{rr}
0 & 1\\
2 & 3
\end{array}\end{bmatrix}\mbox{ and }B=\begin{bmatrix}\begin{array}{rr}
2 & 2\\
2 & 2
\end{array}\end{bmatrix}\mbox{.}
\]
Find $A+B$.

\begin{enumerate}[resume]
\item []$A+B=\begin{bmatrix}\begin{array}{rr}
0+2 & 1+2\\
2+2 & 3+2
\end{array}\end{bmatrix}=\begin{bmatrix}\begin{array}{rr}
2 & 3\\
4 & 5
\end{array}\end{bmatrix}$
\end{enumerate}
\end{enumerate}
\item {[}Bonus{]} Prove that $n\times n$ matrices have an identity element
under matrix addition and multiplication, i.e., that given any $n\times n$
matrix $A$ there exists a matrix $B$ and $C$ such that $A+B=A$
and $AC=A$. 

\begin{proof}
Let 
\[
B=\begin{bmatrix}\begin{array}{rrrr}
0 & 0 & \cdots & 0\\
0 & 0 & \cdots & 0\\
\vdots & \vdots & \ddots & \vdots\\
0 & 0 & \dots & 0
\end{array}\end{bmatrix}\mbox{ and }C=\begin{bmatrix}\begin{array}{rrrr}
1 & 0 & \cdots & 0\\
0 & 1 & \cdots & 0\\
\vdots & \vdots & \ddots & \vdots\\
0 & 0 & \dots & 1
\end{array}\end{bmatrix}
\]
(This is the zero matrix, $\mathbf{0}$, and the identity matrix $I$
for $n\times n$ matrices respectively.) So then, 
\[
A+B=A+\mathbf{0}=\begin{bmatrix}\begin{array}{rrrr}
a_{11}+0 & a_{12}+0 & \cdots & a_{1n}+0\\
a_{21}+0 & a_{22}+0 & \cdots & a_{2n}+0\\
\vdots & \vdots & \ddots & \vdots\\
a_{n1}+0 & a_{n2}+0 & \dots & a_{nn}+0
\end{array}\end{bmatrix}=A
\]

For $AC=AI$ we have that that the entry for the $ij^{\mbox{th}}$
place is $(AC)_{ij}=(a_{i1},a_{i2},\dots,a_{in})\cdot(c_{j1},c_{j2},\dots,c_{jn})$.
But $c_{jk}=0$ for every coordinate except $c_{jj}=1$. So then 
\[
(AC)_{ij}=(a_{i1},a_{i2},\dots,a_{in})\cdot(0,0,\dots,1,\dots,0)=a_{i1}\cdot0+a_{i2}\cdot0+\cdots+1\cdot a_{ij}+\cdots+0\cdot a_{in}=a_{ij}
\]
Thus $AC=AI=A$ (note also that it similarly follows that $\mathbf{0}+A=A$
and $IA=A$; it is usually required that identity elements be commutative.).
\end{proof}
\end{enumerate}

\end{document}
